% Source: https://artofproblemsolving.com/community/c3h2029806_happy_pi_day_maybe_a_competition
% Stable identity: aops-topic-2029806-q07
The square has side \(s=4\).

\emph{Largest equilateral triangle.} Let an equilateral triangle of side \(t\) lie in the square, and let \(\theta\) be the angle between one of its sides and the bottom side of the square. Its width (length of projection) in the horizontal direction is \(t\cos d\), where \(d\in[0^\circ,30^\circ]\) is the distance from \(\theta\) to the nearest multiple of \(60^\circ\) (the projection is the longest side-projection, and the three sides make angles \(\theta,\theta+60^\circ,\theta+120^\circ\)). The vertical width is \(t\cos d'\) with \(d'\) defined the same way for \(\theta+90^\circ\); since \(90^\circ\equiv30^\circ\pmod{60^\circ}\), \(d'=30^\circ-d\). Both widths are at most \(4\), so
\[
4\ge t\max(\cos d,\cos(30^\circ-d))=t\cos(\min(d,30^\circ-d))\ge t\cos15^\circ .
\]
Thus \(t\le\dfrac{4}{\cos15^\circ}=\dfrac{16}{\sqrt6+\sqrt2}=4(\sqrt6-\sqrt2)\). This is attained by the triangle with one vertex at a corner of the square and its two sides from that corner making \(15^\circ\) angles with the square's sides:

\begin{asy}
usepackage("amsmath");
size(4cm);
defaultpen(fontsize(10pt));
real s=4, tt=Tan(15);
pair O=(0,0), P=(s,s*tt), Q=(s*tt,s);
draw((0,0)--(s,0)--(s,s)--(0,s)--cycle,linewidth(0.9));
draw(O--P--Q--cycle,blue+linewidth(0.9));
real a=abs(P-Q); pair I=(O+P+Q)/3; real r=a/(2*sqrt(3));
draw(circle(I,r),red);
label("$15^\circ$",O+1.3*dir(7.5),fontsize(8pt)); label("$15^\circ$",(0.24,2.2),fontsize(7pt));
label("$4$",(s/2,0),S); dot(I); label("$r$",I+(r/2,0),N,red); draw(I--I+(r,0),red);
label("\mbox{}",(min(currentpicture,user=true)+max(currentpicture,user=true))/2,black); // keeps the text after the figure black when LaTeX typesets the labels inline
\end{asy}

its vertices \((0,0)\), \((4,4\tan15^\circ)\), \((4\tan15^\circ,4)\) are pairwise \(4(\sqrt6-\sqrt2)\) apart (using \(\tan15^\circ=2-\sqrt3\)).

\emph{The circle.} The largest circle in a triangle is its incircle; for an equilateral triangle of side \(t\) its radius is \(r=\dfrac{t}{2\sqrt3}\). So
\[
r^2=\dfrac{t^2}{12}=\dfrac{16(8-4\sqrt3)}{12}=\dfrac{32-16\sqrt3}{3},
\]
and the area of the circle is
\[
\pi r^2=\boxed{\dfrac{16(2-\sqrt3)\pi}{3}}\approx1.40 .
\]
